\chapter{Real Physics without Philosophy}\label{ch:philosophy}
\index{philosophy of physics}\index{ontology}\index{measurement problem}\index{Heisenberg cut}\index{interpretation of quantum mechanics}\index{configuration space}

\begin{quote}
\itshape A theory that can be wrong is doing physics; a theory that cannot say what it is about is doing something
else.\\
\normalfont\hfill --- working note, 2026.
\end{quote}
\bigskip

A physical theory is an account of what exists and how it behaves; prediction is derivative --- the behaviour of
what exists, computed forward. By this measure standard quantum mechanics is not, or not yet, real physics: it is
an \emph{epistemology}, a calculus for the outcomes of specific experiments that declines to say what is
physically present between them. This chapter argues that this is not a stance one may take without cost, that the
century of ``interpretations'' is the symptom of a single absence, and that RealQM --- being an ontology of charge
in ordinary space --- removes the need for the interpretive apparatus rather than adding to it.

\section{Physics is ontology; prediction is derivative}\label{sec:physics-ontology}

A physical theory tells us what the world is made of and how it behaves; from that, what we will see follows.
Newtonian mechanics posits masses at positions with momenta, evolving under forces, and its predictions are the
computed trajectories of those real things. Maxwell's theory posits fields in space; continuum mechanics posits
densities, stresses, and velocities. In every case the prediction is \emph{derivative}: the forward computation
of the behaviour of something the theory holds to exist. Remove the ontology and there is nothing whose
behaviour is being computed. Physics without ontology is not a leaner physics; it is, at most, a rule for
anticipating readings --- a different and lesser thing, and, as we now argue, one that cannot stand on its own.

\section{Prediction presupposes ontology: the Heisenberg cut}\label{sec:heisenberg-cut}

It is often held that quantum mechanics is the exception --- that it predicts the outcomes of measurements and
need commit to nothing further. But a prediction is a statement \emph{about something}: to predict an outcome is
to predict that some real thing will be found in some real state. Standard QM's ``outcomes'' are readings of an
apparatus --- a pointer, a click, a track --- which are macroscopic, classical, and taken as real. So the theory
does have an ontology after all: the classical apparatus. What it refuses is an ontology for the microscopic
system, represented only by an amplitude.

The result is the \emph{Heisenberg cut}: a line, nowhere fixed by the theory, above which things are real and
below which they are only amplitudes. Two failures follow. First, the theory cannot say where the cut falls,
because nothing in the formalism marks it. Second, and worse, it grants to the macroscopic the very reality it
denies to the microscopic, although the macroscopic is \emph{made of} the microscopic: the pointer is a swarm of
the same electrons and nuclei the theory says have no definite state. This is not a subtlety but an incoherence.
Prediction without ontology is thus not economical; it is contradictory --- it must help itself, at the
macroscopic end, to exactly the ontology it refuses at the microscopic.

\section{The measurement problem is the missing ontology}\label{sec:measurement}

Seen this way, the measurement problem is not a deep feature awaiting a clever resolution; it is the direct
expression of the missing ontology. Ask \emph{what is measured}, and the theory has no answer, because the object
it calls fundamental --- the amplitude --- is not a thing that could be measured. ``Collapse'' is the name for
the step in which the theory must, without warrant from its own dynamics, convert the amplitude into a definite
fact at the moment a measurement is declared. Schr\"odinger's cat is the reductio of pretending the amplitude was
the reality all along. Each is a symptom of one absence: a theory that said what exists would say what is
measured, and would not need to collapse anything.

\section{One root: the wave function over configuration space}\label{sec:config-space}

The absence has a single source. The fundamental object of standard QM, $\Psi(x_1,\dots,x_N)$, lives on a
$3N$-dimensional \emph{configuration space}, not in the three-dimensional space we and our instruments occupy.
Kohn stated the working consequence in his Nobel lecture: the many-electron wave function is not a legitimate
scientific concept beyond a handful of electrons, being unrecordable even in principle. A function on a space
that is not physical space can carry a probability distribution, but it cannot be an ontology. From this one fact
the familiar troubles descend --- the \emph{arena} problem (in what space does the state exist?), the
\emph{individuality} problem (identical particles with no separate existence), and the nonlocality of a
correlated whole that has no parts in space to be local to. The puzzles of quantum foundations are not many
independent mysteries; they are one, the absent ontology of an object in the wrong space, wearing several faces.

\section{Interpretation as symptom}\label{sec:interpretation}

The century of ``interpretations'' is the diagnostic. Copenhagen declares that the microscopic has no properties
until measured --- it accepts the absence and forbids the question. Many-worlds keeps the amplitude as real and
multiplies the world to fit it. De Broglie--Bohm restores particles with definite positions beneath the
amplitude. Objective-collapse models change the dynamics so the amplitude localises. However they differ, each is
an attempt either to supply the ontology the formalism lacks or to legislate that none is needed. That so many
mutually incompatible interpretations coexist, none forced by the equations, is the mark not of a rich theory but
of an \emph{underdetermined} one. A theory that said what exists would admit no interpretations, for there would
be nothing left to interpret. Philosophy of physics, in its dominant modern form, simply \emph{is} the labour of
coping with a physics that withheld its own ontology.

\section{RealQM: an ontology that predicts real physics}\label{sec:realqm-ontology}

RealQM removes the need at the root. Its fundamental object is not an amplitude on configuration space but $N$
charge densities $\psi_i^2$, each carried on its own region $\Omega_i$ of ordinary three-dimensional space, the
regions non-overlapping. This is an ontology in the plain sense: it says what exists --- extended charge in
space --- and where. It is the same \emph{kind} of object macroscopic continuum mechanics works with; an atom is
a continuum-mechanics problem an \aa ngstr\"om across. From it real physics is computed --- the total energies of
atoms across the periodic table, covalent and hydrogen bonds, and, reading the nucleus as protons and electrons
bound by Coulomb alone, nuclear binding energies --- and what is computed is not the reading of an instrument but
the thing itself.

The foundational puzzles do not arise, because their common source is gone. There is no Heisenberg cut, because
everything is real at every scale. There is no collapse, because the evolution is deterministic and the density
is definite at all times; ``measurement'' is the ordinary interaction of one real charge distribution with
another. The apparent randomness of quantum experiments is the epistemic randomness of initial configurations we
neither see nor control, as in any deterministic account of a complicated system --- ours, not the world's. We do
not claim RealQM is complete: it does not yet deliver the spin-exchange phenomena, and the boundary of its reach
is still being mapped (Chapter~\ref{ch:limits}). But its incompleteness is the \emph{ordinary} incompleteness of
a physical theory --- there is more to compute, and the claims can be checked and found wanting --- not the
structural incoherence of a theory that cannot say what it is about. A theory that can be \emph{wrong} is doing
physics; and where RealQM is right it is right about \emph{things}, not about readings.

\section{The wager, and its reach}\label{sec:wager}

It is worth setting down, plainly and as a conditional, exactly what RealQM claims and how far the claim extends.
The claim is a wager about ontology: that the physics of matter is, at bottom, \emph{real charge in ordinary
three-dimensional space}. If that is so, RealQM --- the dynamics of exactly that --- can capture it. The
syllogism is valid in form; its force rests on two loads.

\emph{First: how much of physics is real charge?} Where the answer is ``essentially all of it'' --- the
Coulomb-structured world of atoms, molecules, bonds, the periodic table, and a growing list of magnetic and
spectroscopic observables --- RealQM captures the physics, often more cleanly than the configuration-space
formalism, because that domain simply \emph{is} the electrostatics of real charge. But some of physics is not
manifestly charge alone: light is the electromagnetic \emph{field}, a second real entity, so radiation lives in a
\emph{charge-plus-field} coupling (Chapter~\ref{ch:sm-coupling}); the statistics of measurement across
incompatible settings are informational as much as material; gravity and the precision relativistic sector stand
outside the present theory. Where the premise is incomplete, so is the reach.

\emph{Second: even granting the ontology, is the dynamics right?} Possessing the right stuff --- charge in real
space --- does not by itself fix the right law. RealQM's particular functional (Coulomb, no self-interaction,
non-overlapping domains, kinetic confinement) is itself a hypothesis, confirmed by its hits (total energies,
bonds, helium's diamagnetic susceptibility) and corrected by its misses (many-electron valence sizes, exchange
pairing). A true ontology still owes a true equation.

So the honest thesis is neither ``RealQM explains everything'' nor ``some residue defeats it,'' but a conditional
under adjudication: \textbf{where physics is real charge in three-dimensional space, RealQM captures it; the open
frontier is to measure how far that reaches, phenomenon by phenomenon} --- extending the picture where charge
alone is not enough, and marking honestly where it meets a limit. That discipline --- capture, extend, or
concede, case by case --- is what makes it physics rather than doctrine.

\section{Real physics without philosophy}\label{sec:no-philosophy}

The conclusion is in the title. Physics with a clear ontology is simply physics: it says what exists, computes
what it does, and is tested against the world. It needs no interpretation, because it has already said what it is
about --- and one does not interpret a bridge calculation or a weather model, one checks it. The elaborate
philosophy that surrounds quantum mechanics --- the interpretations, the measurement debates, the metaphysics of
the wave function --- is not intrinsic to the physics of the small. It is the price of a formalism that left
three-dimensional space and, with it, its ontology. Return the theory to real space, let the charge densities be
the things they plainly appear to be, and the philosophy is not refuted so much as \emph{retired}: there is
nothing left for it to do. The mess was never in nature; it was in the arena. Change the arena back, and physics
is once again about what there is.
